
\include{./def/yf-formatting}
\include{./def/yf-def}

\usepackage{wasysym}
%\usepackage{times}
%\usepackage{mathabx}

\title{}
\author{}
\date{}


%\def\A{\mathcal{A}}
%\def\B{\mathcal{B}}
\def\C{\mathcal{C}}
%\def\D{\mathcal{D}}
\def\E{\mathcal{E}}
\def\FF{\EuScript{F}}
\def\G{\mathcal{G}}
\def\GG{\EuScript{G}}
\def\H{\mathcal{H}}
\def\I{\mathcal{I}}
\def\J{\mathcal{J}}
\def\L{\mathcal{L}}
\def\Q{\mathcal{Q}}
\def\S{\mathcal{S}}
\def\T{\mathcal{T}}
\def\V{\mathcal{V}}
%\def\W{\EuScript{W}}
\def\W{\mathscr{W}}
\def\X{\mathcal{X}}
\def\Y{\mathcal{Y}}
\def\Z{\mathcal{Z}}
\def\att{\mit{diffatt}}
\def\boldatt{\mathbf{att}}
\def\dom{\mathbf{dom}}
\def\ext{\mathit{ext}}
\def\join{\mathit{join}}
\def\schema{\mathit{schema}}
\def\trans{\mathsf{T}}

\def\extraspacing{\vspace{4mm} \noindent}
\def\vgap{\vspace{4mm}}
\def\vslit{\vspace{2mm}}

\begin{document}

\boxminipg{\linewidth}{
    \begin{center}
        \vspace{3mm}
    {\Large
    (CE7456 Lecture Notes 11) Tightness of the AGM Bound} \\[2mm]
    {\large Yufei Tao}
    \vspace{3mm}
    \end{center}
}

Today, we prove the tightness of the AGM bound on joins.

\vgap

Let $\boldatt$ be a finite set, where each element is called an {\em attribute}. For a non-empty set $\X \subseteq \boldatt$ of attributes, a {\em tuple} over $\X$ is a function $\bm{u}: \X \rightarrow \intdom$ (where $\intdom$ represents the set of integers); for each $X \in \X$, we refer to $\bm{u}(X)$ as the {\em value} of $\bm{u}$ on $X$. For any non-empty subset $\Y \subseteq \X$, we define $\bm{u}[\Y]$ --- called the {\em projection} of $\bm{u}$ on $\Y$ --- as the tuple $\bm{v}$ over $\Y$ satisfying $\bm{v}(Y) = \bm{u}(Y)$ for every attribute $Y \in \Y$. A {\em relation} $R$ is a set of tuples over the same set $Z$ of attributes; we refer to $Z$ as the {\em schema} of $R$ and represent it as $\schema(R)$.

\vgap

%Our discussion in this lecture assumes that $\schema(\Q)$ has a constant size.




%\vgap


We define a {\em join} as a set $\Q$ of relations. Let $\schema(\Q) = \bigcup_{R \in Q} \schema(R)$. The result of $\Q$ is a relation over $\schema(\Q)$ formalized as:
\myeqn{
    \join(\Q) &=& \{\textrm{tuple $\bm{u}$ over $\schema(\Q)$} \mid \forall R \in \Q: \bm{u}[\schema(R)] \in R\}. \label{eqn:join-rslt}
}
The {\em schema graph} of $\Q$ is the hypergraph $\H = (\V, \E)$ where
\myeqn{
    \V &=& \schema(\Q) \nn \\
    \E &=& \set{\schema(R) \mid R \in \Q}. \nn
}
Note that there is one-one correspondence between the relations in $\Q$ and the hyperedges in $\E$. For each hyperedge $e \in \E$, let $R_e$ represent the corresponding relation $R \in \Q$.

\vgap

A vector $\W$ over the space $\E$ is a {\em fractional edge cover} of $\H$ if
\myitems{
    \item $\W_e \ge 0$ for every $e \in \E$, and
    \item $\sum_{e \in \E : X \in e} \W_e \ge 1$ for every $X \in \V$.
}
Every such $\W$ is associated with an {\em AGM value} defined as
\myeqn{
    \mit{AGM}_\W &=& \prod_{e \in \E} |R_e|^{\W_e}. \label{eqn:agm}
}
The {\em AGM bound} of $\Q$ is $\min_\W \mit{AGM}_\W$ where the minimization is taken over all fractional edge covers $\W$ of $\H$.

\vgap


Our purpose is to prove:

\begin{theorem} \label{thm:agm-tight}
   Fix an arbitrary hypergraph $\H = (\V, \E)$ where $|\V|$ is a constant, and an arbitrary integer vector $N$ over $\E$ such that $N_e \ge 1$ for each $e \in \E$. There exists a join $\Q$ satisfying
   \myitems{
        \item the schema graph of $\Q$ is $\H$;
        \item for each $e \in \E$, $|R_e| \le N_e$;
        \item $|\join(\Q)|$ is at least the AGM bound of $\Q$ up to a constant factor.
   }
\end{theorem}

\extraspacing {\bf Math Conventions.} For an integer $x \ge 1$, the notation $[x]$ denotes the set $\{1, 2, ..., x\}$. Every logarithm $\log(\cdot)$ has base $2$.

\section{Linear Programming} \label{app:lp}

The material of this section can be found in most textbooks (e.g., \cite{mg07}) on linear programming (LP) and is included for self-containment reasons.

\vgap

Suppose that $\mb{A}$ is an $m \times n$ matrix, $\bm{c}$ is an $n \times 1$ matrix (a.k.a.\ an $n$-dimensional vector), and $\bm{b}$ an $m \times 1$ matrix (a.k.a.
an $m$-dimensional vector). Consider an LP of the form:
\myeqn{
    \text{find an $n \times 1$ matrix $\bm{x}$ to maximize $\bm{c}^\trans \bm{x}$ subject to $\mb{A}\bm{x} \le \bm{b}$ and $\bm{x} \ge 0$}.
    \label{eqn:lp:primal}
}
The duality of the above LP has the form:
\myeqn{
    \text{find an $m \times 1$ matrix $\bm{y}$ to minimize $\bm{b}^\trans \bm{y}$ subject to $\mb{A}^\trans\bm{y} \ge \bm{c}$ and $\bm{y} \ge 0$.}
    \label{eqn:lp:dual}
}
We will refer to \eqref{eqn:lp:primal} as the {\em primal form} and to \eqref{eqn:lp:dual} as the {\em dual form}.

\vgap

A {\em feasible solution} to the LP of the form \eqref{eqn:lp:primal} is an $n \times 1$ matrix $\bm{x}$ satisfying $\mb{A}\bm{x} \le \bm{b}$ and $\bm{x} \ge 0$. The result of the LP has three possible cases:
\myitems{
    \item {{\em empty}}: no feasible solutions exist;
    \item {{\em unbounded}}: for any finite real value $\tau$, there is a feasible solution $\bm{x}$ such that $\bm{c}^\trans \bm{x} \ge \tau$;

    \item {{\em bounded}}: $\bm{c}^\trans \bm{x}$ takes a finite maximum --- called the {\em optimal value} --- among all feasible solutions $\bm{x}$.
}
In the bounded case, we define an {\em optimal solution} to the LP as a feasible solution $\bm{x}^*$ such that $\bm{c}^\trans \bm{x}^*$ equals the optimal value.

\vgap

Likewise, a {\em feasible solution} to the LP of the form \eqref{eqn:lp:dual} is an $m \times 1$ matrix $\bm{y}$ such that $\mb{A}^\trans \bm{y} \ge \bm{c}$ and $\bm{y} \ge 0$. The result of the LP has three possible cases:
\myitems{
    \item {{\em empty}}: no feasible solutions exist;
    \item {{\em unbounded}}: for any finite real value $\tau$, there is a feasible solution $\bm{y}$ such that $\bm{b}^\trans \bm{y} \le \tau$;

    \item {{\em bounded}}: $\bm{b}^\trans \bm{y}$ takes a finite minimum --- called the {\em optimal value} --- among all feasible solutions $\bm{y}$.
}
In the bounded case, we define an {\em optimal solution} to the LP as an $m \times 1$ matrix $\bm{y}^*$ such that $\bm{b}^\trans \bm{y}^*$ equals the optimal value.

\vgap

The {\em strong duality theorem} states:
\myitems{
    \item if the primal form is bounded, the dual form is bounded and returns the same optimal value;

    \item if the dual form is bounded, the primal form is bounded and returns the same optimal value.
}


%An {\em optimal solution} to the LP in \eqref{eqn:lp:primal} is an $n \times 1$ matrix (a.k.a.\ vector) $\bm{x}^*$ such that $\bm{c}^\trans \bm{x}^*$ equals the optimal value. Similarly, an {\em optimal solution} to the LP in \eqref{eqn:lp:dual} is an $m \times 1$ matrix (a.k.a.\ vector) $\bm{y}^*$ such that $\bm{b}^\trans \bm{y}^*$ equals the optimal value. By the strong duality theorem, $\bm{c}^\trans \bm{x}^* = \bm{b}^\trans \bm{y}^*$.

\section{Proof of Theorem~\ref{thm:agm-tight}} \label{sec:proof}

Finding a fractional edge cover $\W$ to minimize \eqref{eqn:agm} can be cast as an LP problem:
\minipg{0.8\linewidth}{
    {\bf minimize} $\sum_{e \in \E} \W_e \log |R_e|$ {\bf subject to}
    \myeqn{
        \W_e \ge 0 &&  \forall e \in \E \nn \\
        \sum_{e \in \E: X \in e} \W_e \ge 1 &&  \forall X \in \V. \nn
    }
}
The result of the LP (i) cannot be unbounded because clearly $\sum_{e \in \E} \W_e \log |R_e| \ge 0$, and (ii) cannot be empty because setting $\W_e = 1$ for all $e \in \E$ yields a feasible solution. Therefore, the LP must have a (finite) optimal value.

\vgap

Motivated by the observation, we formulate two LP instances --- using the hypergraph $\H = (\V, \E)$ and vector $N$ (over $\E$) given in the statement of Theorem~\ref{thm:agm-tight} --- that are dual to each other. Set
\myeqn{
    n &=& |\V| \nn \\
    m &=& |\E| \nn
}
Let us list the vertices in $\V$ as $X_1, X_2, ..., X_n$ and the edges in $\E$ as $e_1, e_2, ..., e_m$ (the vertex and edge orderings are not important).

\vgap

Define an $m \times n$ matrix $\mb{A}$ where for each $i \in [m]$ and $j \in [n]$:
\myeqn{
    \mb{A}[i,j] &=& \left\{
        \begin{tabular}{ll}
            1 & if $e_i$ contains $X_j$ \\
            0 & otherwise
        \end{tabular}
    \right. \nn
}
Define an $m \times 1$ vector $\bm{b}$ where for each $i \in [m]$:
\myeqn{
    \bm{b}[i] = \log N_{e_i} \nn
}
and an $n \times 1$ vector $\bm{c}$ where for each $j \in [n]$:
\myeqn{
    \bm{c}[j] = 1. \nn
}

\vgap

The above formulation yields a primal LP of the form \eqref{eqn:lp:primal} and its dual LP of the form \eqref{eqn:lp:dual}. The primal LP can also be written as:
\minipg{0.8\linewidth}{
    {\bf maximize} $\sum_{j=1}^n \bm{x}[j]$ {\bf subject to}
    \myeqn{
        \bm{x}[j] \ge 0 &&  \forall j \in [n] \nn \\
        \sum_{j \in [n]: X_j \in e_i} \bm{x}[j] \le \log N_{e_i} &&  \forall i \in [m]. \nn
    }
}
The dual LP can be written as:
\minipg{0.8\linewidth}{
    {\bf minimize} $\sum_{i=1}^m \bm{y}[i] \cdot \log N_{e_i}$ {\bf subject to}
    \myeqn{
        \bm{y}[i] \ge 0 &&  \forall i \in [m] \nn \\
        \sum_{i \in [m]: X_j \in e_i} \bm{y}[i] \ge 1 &&  \forall j \in [n]. \nn
    }
}
The results of both LPs must be bounded. Denote by $\bm{x}^*$ (resp., $\bm{y}^*$) an optimal solution to the primal (resp., dual) LP. By the strong duality theorem, we have:
\myeqn{
    \sum_{j=1}^n \bm{x}^*[j]
    &=&
    \sum_{i=1}^m \bm{y}^*[i] \cdot \log N_{e_i}.
    \label{eqn:lp-opt}
}

\vgap

For any join $\Q$ such that (i) the schema graph of $\Q$ is $\H$ and (ii) $|R_e| \le N_e$ for each $e \in \E$, the AGM bound of $\Q$ is equivalent to
\myeqn{
    2^
    {\sum_{i=1}^m \bm{y}^*[i] \cdot \log N_{e_i}}.
    \label{eqn:proof:agm}
}
Next, we construct such a join $\Q$ whose result contains
\myeqn{
    2^
    {\sum_{j=1}^n \bm{x}^*[j]}
    \label{eqn:proof:join-size}
}
up to a constant factor. By \eqref{eqn:lp-opt}, the values in \eqref{eqn:proof:agm} and \eqref{eqn:proof:join-size} are identical. This will then complete the proof of Theorem~\ref{thm:agm-tight}.

\vgap

For each $j \in [n]$, define the {\em domain} of vertex $X_j$ as
\myeqn{
    \dom(X_j) &=& \set{1, 2, ..., 2^{\lf \bm{x}^*[j] \rf}}.
    \nn
}
For each $i \in [m]$, define the relation $R_{e_i} \in \Q$ --- namely, the relation whose schema corresponds to $e_i$ --- as the cartesian product of the domains of the vertices in $e_i$. Formally, assume that $e_i$ has vertices $X_{j_1}, X_{j_2}, ..., X_{j_t}$ for some $t \ge 1$. Then, $R_{e_i}$ has
\myeqn{
    \prod_{k=1}^t |\dom(X_{j_k})| \nn
}
tuples such that, for any $(x_{j_1}, x_{j_2}, ..., x_{j_t}) \in \dom(X_{j_1}) \times \dom(X_{j_2}) \times ... \times \dom(X_{j_t})$, the relation $R_{e_i}$ contains a tuple $\bm{u}$ with $\bm{u}(X_{j_k}) = x_{j_k}$ for all $k \in [t]$.

\vgap

By our construction, for any $(x_1, x_2, ..., x_n) \in \dom(X_1) \times \dom(X_2) \times ... \times \dom(X_n)$, the result of the join $\Q$ must contain a tuple $\bm{u}$ satisfying $\bm{u}(X_j) = x_j$ for each $j \in [n]$. The size of $\join(\Q)$ is
\myeqn{
    \prod_{j=1}^n 2^{\lf \bm{x}^*[j] \rf}
    \ge
    \fr{1}{2^n} \cdot \prod_{j=1}^n 2^{\bm{x}^*[j]}
    =
    \fr{1}{2^n}
    \cdot
    2^
    {\sum_{j=1}^n \bm{x}^*[j]}.
    \nn
}
Since $n$ is a constant, the above can be lower than \eqref{eqn:proof:join-size} by at most a constant factor.

\section{Remark}

The proof in Section~\ref{sec:proof} is taken from \cite{s23}.

\bibliographystyle{abbrv}
\bibliography{ref}

\end{document}
